\section{The precise constructive setting}
\label{sec:foundations}

\subsection{The statements being proved}

The linked Hurwitz theorem concerns infinitely many reduced fractions
$p/q$ satisfying $|x-p/q|<1/(\sqrt5 q^2)$ for each irrational $x$.
The linked Khinchin theorem concerns almost-everywhere approximation by
$|x-p/q|<\psi(q)/q$ and the series $\sum_q\psi(q)$.
Their literal classical formulations quantify over real numbers, infinite
sequences, and Lebesgue measurable sets.  First-order $\ipra$ has one
numerical sort, primitive recursive function symbols, intuitionistic
first-order logic with equality, and induction only for quantifier-free
formulas.  We therefore specify the finite arithmetic content of the
constructive versions before using real-number language.

The two central output assertions are as follows.
\begin{enumerate}
\item \textbf{Hurwitz witness.} For a fixed primitive recursive rational
name $x_s$ with error at most $2^{-s}$, a fixed primitive recursive
rational-apartness modulus, and every natural $B$, compute $p,q,s$ with
$q>B$, $\gcd(p,q)=1$, and
\[
 5q^4\bigl(|x_s-p/q|+2^{-s}\bigr)^2<1.
\]
This is a decidable inequality between rational numbers.  Its proof gives
the strict classical coefficient, rather than a coefficient enlarged by
an arbitrary epsilon.
\item \textbf{Khinchin finite coverage.} For fixed nonnegative rational
primitive recursive $\psi$ and an appropriate primitive recursive modulus,
compute finite rational interval lists.  In the divergent case the
output $U(N,m)$ is labelled by reduced fractions with denominators
between $N$ and a computed $M(N,m)$ and satisfies
\[
 U(N,m)\subseteq\bigcup_{q=N}^{M(N,m)}
       \bigcup_{p=0}^q (p/q-\psi(q)/q,p/q+\psi(q)/q),
 \qquad |[0,1]\setminus U(N,m)|<2^{-m}.
\]
In the convergent case every finite portion of a specified tail has
length at most the specified tolerance.  The measure symbols here mean
finite rational length computed by endpoint sorting.
\end{enumerate}

The divergence proof needs only that $\psi$ be nonincreasing.  The linked
encyclopedia page states the stronger older hypothesis that $q\psi(q)$
is nonincreasing, in the same normalization $\psi(q)/q$.  That stronger
hypothesis implies ours: for $q\le r$, $r\psi(r)\le q\psi(q)$ implies
$\psi(r)\le(q/r)\psi(q)\le\psi(q)$.  The present proof consequently
covers the linked formulation as well as the usual weaker monotonicity
assumption.  The conversion to the alternative radius notation
$\Psi(q)=\psi(q)/q$ is stated separately below.

\subsection{Data, hypotheses, and theorem schemata}

``Primitive recursive'' always means obtained from zero, successor and
projections by composition and primitive recursion.  Natural numbers
code signed integers, fractions with positive denominator, finite lists,
and finite machine states.  Numerator signs are stored explicitly.
Rational arithmetic is implemented by integer arithmetic and comparison
by cross multiplication.  Reduction of fractions is optional except where
coprimality is required; gcd can be computed by the bounded search through
divisors up to the larger input.  List decoding and sorting can be performed
by bounded loops on the code and its recorded length.  Efficiency is not
needed for any foundational assertion here.

An input function such as $\psi$, $x_s$, or a modulus is a \emph{fixed}
primitive recursive definition, represented by its own function symbol.
The statements are schemata, one for every such fixed collection of
definitions.  Their hypotheses need not already be theorems of $\ipra$:
the formal statement is the implication from the displayed coherence,
monotonicity, separation, or mass assumptions.  For a particular example,
one additionally proves those assumptions.  Alternatively all queried
values and their finite verifications may be supplied as finite tables.
There is no first-order quantification over arbitrary real functions and
no supposed primitive recursive universal evaluator for all primitive
recursive program codes.

A name $s\mapsto x_s\in\Q$ may be specified arithmetically by
\[
 |x_s-x_t|\le2^{-s}+2^{-t}\qquad(s,t\in\N).
\]
The notation $|x-x_s|\le2^{-s}$ expresses its ordinary interpretation;
the internal arguments use rational enclosures and these coherence
inequalities.  A positive separation modulus supplies quantitative
apartness from every rational with bounded denominator.  Mere negative
irrationality, expressed as the absence of an equality, is not being
silently promoted to that modulus.

Similarly we do not decide whether an arbitrary nonnegative series
converges or diverges.  We prove separate implications from a supplied
tail bound or a supplied divergence witness.  No procedure extracting a
primitive recursive modulus from an unqualified classical convergence or
divergence assertion is part of either theorem.  Section~\ref{sec:extensions}
extends the metric argument to real-valued approximation functions with
primitive recursive names and witnessed tail information.

\subsection{Bounded predicates and quantifier-free induction}

Every bounded predicate on the finite data used here has a primitive
recursive Boolean characteristic function.  For a Boolean predicate
$b(i)$, conjunction over $i<n$ is the finite product of the bits, and
existence is a bounded search with a failure flag.  Nesting these
constructions handles finitely nested bounded quantifiers.

There is a relevant proof-theoretic detail.  We do not apply the iPRA
induction rule directly to a formula containing a bounded quantifier.
We first replace its bounded tests by the corresponding characteristic
terms, prove their specifications, and induct on the resulting equation.
For example the equations
\[
 a(0)=1,\qquad a(n+1)=a(n)b(n)
\]
are accompanied by a bounded first-failure function.  Quantifier-free
induction proves that a failure flag identifies an index $i<n$ with
$b(i)=0$, and that $a(n)=1$ forces $b(i)=1$ for each given $i<n$.
The converse follows from the failure witness and the decidable bit
case distinction.  This proves the bounded-universal specification
without induction on a quantified formula.  The same method handles
finite maxima, bounded searches, and endpoint membership tests.

More generally, let $S(0,z)=G(z)$ and
$S(n+1,z)=F(n,S(n,z),z)$ be a fixed primitive recursive state iteration.
If $V(n,s,z)$ is a primitive recursive Boolean invariant and its initial
and preservation equations have been proved, iPRA's quantifier-free
induction proves $V(n,S(n,z),z)=1$.  Universal assumptions about the input
names or moduli may remain in the context and be instantiated at each
step; they are not induction formulas.  The invariants used below are
listed explicitly in Section~\ref{sec:formalization}.  Thus the reasoning
is a finite arithmetic proof with a specified induction discipline, not
an appeal to the rule that every terminating algorithm has iPRA-provable
correctness.  That rule would be unjustified.

\subsection{Finite length and finite second moments}

For rational open intervals, retain their endpoints and their open/closed
flags when membership matters.  Sort the finite list of all endpoints.
Between successive distinct endpoints, membership in each input interval
is constant and can be evaluated at the rational midpoint.  The length
of a union is the sum of the lengths of exactly those cells whose
membership bit is one.  Single endpoints have length zero, but a singleton
in a complement must still be retained for later open-cover statements.
This procedure proves finite additivity on disjoint cells and finite
subadditivity by the inequality between membership bits.

If intervals only touch, their open union omits the common endpoint.
It is permissible to merge strict overlaps; merging touching intervals
would lose that singleton.  The algorithms and proof below retain it.
The two endpoints of $[0,1]$ may safely be included as extra singleton
gaps, even if a relatively open interval already covers them.

For a nonnegative integer-valued function $F$ constant on these cells,
the notation $\int F$ abbreviates a finite rational sum
$\sum_i w_iF_i$.  In particular the only Cauchy--Schwarz inequality needed
in the metric proof follows from
\[
 \left(\sum_i w_i\right)\left(\sum_i w_i v_i^2\right)
   -\left(\sum_i w_i v_i\right)^2
 =\sum_{i<j}w_iw_j(v_i-v_j)^2\ge0.
\]
This is a polynomial identity and a sum of nonnegative rational terms.
It uses no integration theory, completeness theorem, or countable choice.
All appearances of measure inside the main proof are of this form.

\subsection{Relation to the classical literature}

The sharp Hurwitz theorem and its three-convergent proof are classical;
an accessible author-written account is Filaseta's Theorem~25
\cite{filaseta}.  The organization of the metric proof belongs to the
local-distribution and overlap methods underlying local ubiquity
\cite{bv,bdv}.  We prove the counting, overlap, and amplification estimates
in full, rather than importing the general measure-theoretic theorem.
The specific constants are chosen to simplify verification.  We make no
claim that these constants are optimal, or that the classical mathematical
ideas are new.
